\chapter{Data Understanding II -- Multivariate Analysis, Correlation, and Outlier Detection}
\label{chap:dm-data-understanding-multivariate-outliers}

Multivariate analysis studies simultaneous relationships among two or more attributes within a dataset, overcoming the limitations of isolated univariate descriptions. In this chapter, we explore multidimensional visualization techniques, formal indices for linear and non-linear association and correlation, methodological criteria for outlier diagnosis, and an operational diagnostic checklist to conclude the Data Understanding phase.

The case study used as our benchmark is the famous \textbf{Iris} dataset (introduced by botanist Edgar Anderson and popularized by Ronald Fisher in 1936), composed of $150$ observations evenly distributed across $3$ balanced classes of $50$ samples each (\textit{Iris setosa}, \textit{Iris versicolor}, \textit{Iris virginica}) and described by $4$ continuous features measured in centimeters:
\begin{itemize}[noitemsep]
    \item Sepal length (\textit{Sepal Length})
    \item Sepal width (\textit{Sepal Width})
    \item Petal length (\textit{Petal Length})
    \item Petal width (\textit{Petal Width})
\end{itemize}

% ====================================================================
% SECTION 3.1: MULTIVARIATE VISUALIZATION TECHNIQUES
% ====================================================================
\section{Multidimensional Visualization Techniques}
\label{sec:dm-multivariate-visualization}

When exploratory analysis extends to bivariate and multivariate configurations, graphical representations must provide immediate visual insight into geometric class separability, natural cluster formations, and structural anomalies.

\subsection{Scatter Plot and Scatter Matrix}

The \textbf{Scatter Plot} pairs two continuous attributes $(X, Y)$, projecting observations as points with Cartesian coordinates $(x_i, y_i) \in \mathbb{R}^2$.
\begin{itemize}
    \item \textbf{Key Capabilities}: Visually highlights clusters, direct or inverse correlation trends, and bivariate outliers that remain undetectable in univariate histograms.
    \item \textbf{Color Encoding}: By assigning distinct colors or markers according to class labels, practitioners can assess the geometric separability induced by the feature pair.
    \item \textbf{Evidence from Iris}: Crossing \textit{Petal Length} and \textit{Petal Width} reveals sharp class separation: \textit{Setosa} forms an isolated compact cluster, whereas \textit{Versicolor} and \textit{Virginica} exhibit only modest boundary overlap (Figure~\ref{fig:dm-iris-scatter}). In contrast, the plane formed by \textit{Sepal Length} and \textit{Sepal Width} shows substantial inter-class confusion between \textit{Versicolor} and \textit{Virginica}.
\end{itemize}

\begin{figure}[htbp]
    \centering
    \includegraphics[width=0.88\textwidth]{iris_scatterplot_real.pdf}
    \caption{Bivariate scatter plots on the Iris dataset: comparing sepal and petal class separability.}
    \label{fig:dm-iris-scatter}
\end{figure}

When the number of numeric features $p$ is moderate ($3 \le p \le 10$), a \textbf{Scatter Matrix (Scatter Plot Matrix - SPLOM)} is employed: a symmetric $p \times p$ grid containing all $\frac{p(p-1)}{2}$ distinct pairwise scatter plots.
Univariate histograms or kernel density estimates (KDE) are typically placed along the main diagonal, enabling simultaneous inspection of univariate distributions and pairwise correlations (Figure~\ref{fig:dm-iris-splom}).

\begin{figure}[htbp]
    \centering
    \includegraphics[width=0.88\textwidth]{iris_splom_real.pdf}
    \caption{Complete Scatter Plot Matrix (SPLOM) across the four attributes of the Iris dataset.}
    \label{fig:dm-iris-splom}
\end{figure}

The \textbf{Three-Dimensional (3D) Scatter Plot} extends Cartesian projection to three coordinates simultaneously $(X_1, X_2, X_3)$. It allows identifying 3D outliers that appear completely normal in 2D projections, though planar displays suffer from visual occlusion and viewpoint dependence.

\subsection{Visualizing High-Dimensional Spaces}

In spaces where dimension count exceeds $p > 3$, standard orthogonal Cartesian projections fail to provide comprehensive views. Alternative coordinate geometries are required.

\subsubsection{Parallel Coordinates}
In parallel coordinates, coordinate axes are not mutually orthogonal; instead, they are drawn as $p$ equidistant, parallel vertical lines, each calibrated to the value range of its corresponding feature.

\begin{itemize}
    \item \textbf{Geometric Mapping}: A $p$-dimensional observation $\mathbf{x}_i = (x_{i1}, x_{i2}, \dots, x_{ip}) \in \mathbb{R}^p$ is represented not as a point, but as a continuous \textbf{polyline} crossing all parallel axes, intersecting the $j$-th vertical axis at height $x_{ij}$.
    \item \textbf{Diagnostic Traits}: Records sharing identical class labels tend to form compact bundles of polylines displaying coherent trajectories. An isolated polyline with an anomalous trajectory immediately flags a potential multivariate outlier.
    \item \textbf{Axis Ordering}: Visual legibility depends strongly on horizontal axis ordering. Placing strongly correlated attributes on adjacent axes reveals parallel bundles (positive correlation) or pencil-shaped X-crossings (negative correlation) (Figure~\ref{fig:dm-iris-parallel}).
\end{itemize}

\begin{figure}[htbp]
    \centering
    \includegraphics[width=0.88\textwidth]{iris_parallel_coordinates.pdf}
    \caption{Parallel coordinates for the Iris dataset with class-stratified trajectories.}
    \label{fig:dm-iris-parallel}
\end{figure}

\subsubsection{Radar Plot (Spider or Star Chart)}
The \textbf{Radar Plot} shares the structural principle of parallel coordinates, but dimensional axes are arranged \textbf{radially emanating from a shared central pole}, separated by equi-spaced angles:
\begin{equation}
    \theta_j = \frac{2\pi (j-1)}{p}, \quad j = 1, \dots, p
\end{equation}

For each data record, feature values determine the vertex distance along the corresponding radial spoke; connecting consecutive vertices produces a closed polygon.
When radar plots are displayed as small individual multiples for each observation across a grid, they are called \textbf{Star Plots}, facilitating rapid geometric comparisons across multivariate instances (Figure~\ref{fig:dm-iris-radar}).

\begin{figure}[htbp]
    \centering
    \includegraphics[width=0.82\textwidth]{iris_radar_plot.pdf}
    \caption{Class-average profiles for the Iris dataset rendered via Radar Plot.}
    \label{fig:dm-iris-radar}
\end{figure}

% ====================================================================
% SECTION 3.2: COVARIANCE AND CORRELATION
% ====================================================================
\section{Bivariate Dependence Analysis: Covariance and Correlation}
\label{sec:dm-covariance-correlation}

Quantitative dependence analysis investigates the strength and direction with which two random variables co-vary.

\subsection{From Variance to Covariance}

While variance measures the dispersion of a single variable around its own expectation, \textbf{Sample Covariance} measures the joint tendency of two numerical variables $X$ and $Y$ to deviate synchronously from their sample means:
\begin{equation}
    \text{Cov}(X, Y) = s_{xy} = \frac{1}{n-1} \sum_{i=1}^n (x_i - \bar{x})(y_i - \bar{y})
\end{equation}

In the special case where $X = Y$, covariance equals sample variance:
\begin{equation}
    \text{Cov}(X, X) = \frac{1}{n-1} \sum_{i=1}^n (x_i - \bar{x})^2 = s_x^2
\end{equation}

\subsubsection{Covariance Matrix}
For a set of $p$ continuous attributes, joint variability is completely captured by the \textbf{Covariance Matrix} $V \in \mathbb{R}^{p \times p}$:
\begin{equation}
    V = \begin{bmatrix}
        \text{Cov}(X_1, X_1) & \text{Cov}(X_1, X_2) & \cdots & \text{Cov}(X_1, X_p) \\
        \text{Cov}(X_2, X_1) & \text{Cov}(X_2, X_2) & \cdots & \text{Cov}(X_2, X_p) \\
        \vdots & \vdots & \ddots & \vdots \\
        \text{Cov}(X_p, X_1) & \text{Cov}(X_p, X_2) & \cdots & \text{Cov}(X_p, X_p)
    \end{bmatrix}
\end{equation}
\begin{itemize}[noitemsep]
    \item Diagonal entries $V_{jj} = \text{Var}(X_j) = s_j^2$ represent marginal sample variances.
    \item Off-diagonal entries $V_{jk} = \text{Cov}(X_j, X_k)$ represent pairwise covariances.
    \item Since $\text{Cov}(X_j, X_k) = \text{Cov}(X_k, X_j)$, matrix $V$ is \textbf{symmetric} ($V = V^T$) and positive semi-definite ($V \succeq 0$).
\end{itemize}

\subsubsection{Geometric Interpretation of Deviation Products}
The sign of sample covariance is governed by the quadrants formed around the bivariate mean $(\bar{x}, \bar{y})$ in Cartesian space:

\begin{figure}[htbp]
    \centering
    \begin{tikzpicture}[scale=0.95]
        \draw[thick, -Stealth, color=gray!80!black] (-3.8, 0) -- (4.2, 0) node[right, font=\small] {$X$};
        \draw[thick, -Stealth, color=gray!80!black] (0, -3.2) -- (0, 3.5) node[above, font=\small] {$Y$};
        
        \draw[dashed, thick, color=navy!70] (0.5, -3.0) -- (0.5, 3.2);
        \draw[dashed, thick, color=navy!70] (-3.5, 0.4) -- (3.8, 0.4);
        
        \node[below right, font=\footnotesize\bfseries, color=navy] at (0.5, 0) {$\bar{x}$};
        \node[above left, font=\footnotesize\bfseries, color=navy] at (0, 0.4) {$\bar{y}$};
        
        % Quadrant I
        \node[draw=navy!30, fill=navy!4, rounded corners=3pt, align=center, font=\scriptsize, inner sep=4pt] at (2.5, 2.6) {
            \textbf{Quadrant I}\\
            $(x_i > \bar{x}) \land (y_i > \bar{y})$\\
            Product: $(+) \cdot (+) = \mathbf{(+)}$
        };
        \fill[crimson] (1.2, 0.9) circle (2.2pt);
        \fill[crimson] (1.7, 1.3) circle (2.2pt);
        \fill[crimson] (2.2, 1.7) circle (2.2pt);
        
        % Quadrant II
        \node[draw=navy!30, fill=navy!4, rounded corners=3pt, align=center, font=\scriptsize, inner sep=4pt] at (-2.0, 2.6) {
            \textbf{Quadrant II}\\
            $(x_i < \bar{x}) \land (y_i > \bar{y})$\\
            Product: $(-) \cdot (+) = \mathbf{(-)}$
        };
        \fill[darkgreen] (-1.2, 1.2) circle (2.2pt);
        
        % Quadrant III
        \node[draw=navy!30, fill=navy!4, rounded corners=3pt, align=center, font=\scriptsize, inner sep=4pt] at (-2.0, -2.4) {
            \textbf{Quadrant III}\\
            $(x_i < \bar{x}) \land (y_i < \bar{y})$\\
            Product: $(-) \cdot (-) = \mathbf{(+)}$
        };
        \fill[crimson] (-0.8, -0.4) circle (2.2pt);
        \fill[crimson] (-1.4, -0.9) circle (2.2pt);
        \fill[crimson] (-2.1, -1.3) circle (2.2pt);
        
        % Quadrant IV
        \node[draw=navy!30, fill=navy!4, rounded corners=3pt, align=center, font=\scriptsize, inner sep=4pt] at (2.5, -2.4) {
            \textbf{Quadrant IV}\\
            $(x_i > \bar{x}) \land (y_i < \bar{y})$\\
            Product: $(+) \cdot (-) = \mathbf{(-)}$
        };
        \fill[darkgreen] (1.6, -0.9) circle (2.2pt);
    \end{tikzpicture}
    \caption{Geometric decomposition of bivariate deviation products across four mean-centered quadrants.}
    \label{fig:dm-covariance-quadrants}
\end{figure}

\begin{itemize}
    \item Points in Quadrants I and III contribute positive products (concordant deviations) $\implies \text{Cov}(X, Y) > 0$.
    \item Points in Quadrants II and IV contribute negative products (discordant deviations) $\implies \text{Cov}(X, Y) < 0$.
    \item Uniformly dispersed observations cancel out across quadrants, yielding $\text{Cov}(X, Y) \approx 0$.
\end{itemize}

\subsubsection{Structural Limitation: Scale Dependence}
Covariance lacks scale invariance: its magnitude scales linearly with measurement units. Converting income $X$ from euros to thousands of euros ($X' = X / 1000$) shrinks covariance by a factor of $1000$, although the underlying statistical link remains identical. Hence, covariance cannot directly compare pairs with distinct units or scales.

\subsection{Pearson Linear Correlation Coefficient}

The \textbf{Pearson Linear Correlation Coefficient ($r_{xy}$)} standardizes covariance by dividing by the product of both sample standard deviations:
\begin{equation}
    r_{xy} = \frac{\text{Cov}(X, Y)}{s_x s_y} = \frac{\sum_{i=1}^n (x_i - \bar{x})(y_i - \bar{y})}{\sqrt{\sum_{i=1}^n (x_i - \bar{x})^2} \sqrt{\sum_{i=1}^n (y_i - \bar{y})^2}}
\end{equation}

\subsubsection{Mathematical Properties of $r_{xy}$}
\begin{enumerate}
    \item \textbf{Bounded Range}: By the Cauchy-Schwarz inequality on centered vectors, the coefficient is strictly bounded within:
    \begin{equation}
        -1 \le r_{xy} \le +1
    \end{equation}
    \item \textbf{Scale Invariance}: It is a dimensionless index invariant under positive affine transformations ($X' = aX + b, Y' = cY + d$ with $a, c > 0$).
    \item \textbf{Perfect Linear Relationship}:
    \begin{itemize}
        \item $r_{xy} = +1$: perfect positive linear correlation; all points lie exactly on a line with positive slope ($y = ax + b$, $a > 0$).
        \item $r_{xy} = -1$: perfect negative linear correlation ($a < 0$).
    \end{itemize}
    \item \textbf{Absence of Linear Relationship}: A value of $r_{xy} \approx 0$ signifies \textbf{absence of linear association}, but does not imply statistical independence. Non-linear deterministic functions (e.g., $Y = X^2$ with $X$ symmetrically distributed around zero) yield $r_{xy} = 0$, despite full underlying determinism.
\end{enumerate}

\begin{figure}[htbp]
    \centering
    \begin{tikzpicture}[scale=0.88]
        % Panel 1: r = +1.0
        \begin{scope}[shift={(0,0)}]
            \draw[->] (0,0) -- (2.5,0) node[below] {$X$};
            \draw[->] (0,0) -- (0,2.3) node[left] {$Y$};
            \draw[thick, color=crimson] (0.3,0.3) -- (2.2,2.0);
            \fill (0.5,0.48) circle (1.5pt);
            \fill (1.0,0.92) circle (1.5pt);
            \fill (1.5,1.38) circle (1.5pt);
            \fill (2.0,1.82) circle (1.5pt);
            \node[below, font=\scriptsize] at (1.25,-0.2) {$r = +1.0$};
            \node[below, font=\tiny, align=center] at (1.25,-0.6) {Perfect Positive\\Linear};
        \end{scope}

        % Panel 2: r = -0.8
        \begin{scope}[shift={(3.8,0)}]
            \draw[->] (0,0) -- (2.5,0) node[below] {$X$};
            \draw[->] (0,0) -- (0,2.3) node[left] {$Y$};
            \draw[thick, dashed, color=navy] (0.3,2.0) -- (2.2,0.4);
            \fill (0.4,1.8) circle (1.5pt);
            \fill (0.7,1.7) circle (1.5pt);
            \fill (1.0,1.2) circle (1.5pt);
            \fill (1.3,1.4) circle (1.5pt);
            \fill (1.6,0.9) circle (1.5pt);
            \fill (1.9,0.5) circle (1.5pt);
            \node[below, font=\scriptsize] at (1.25,-0.2) {$r = -0.80$};
            \node[below, font=\tiny, align=center] at (1.25,-0.6) {Strong Negative\\Linear Association};
        \end{scope}

        % Panel 3: r = 0.0
        \begin{scope}[shift={(7.6,0)}]
            \draw[->] (0,0) -- (2.5,0) node[below] {$X$};
            \draw[->] (0,0) -- (0,2.3) node[left] {$Y$};
            \fill (0.5,0.5) circle (1.5pt); \fill (0.7,1.8) circle (1.5pt);
            \fill (1.2,1.1) circle (1.5pt); \fill (1.4,0.4) circle (1.5pt);
            \fill (1.8,1.7) circle (1.5pt); \fill (2.0,0.8) circle (1.5pt);
            \node[below, font=\scriptsize] at (1.25,-0.2) {$r \approx 0.0$};
            \node[below, font=\tiny, align=center] at (1.25,-0.6) {No Linear\\Pattern};
        \end{scope}

        % Panel 4: r = 0.0 Non-linear
        \begin{scope}[shift={(11.4,0)}]
            \draw[->] (0,0) -- (2.5,0) node[below] {$X$};
            \draw[->] (0,0) -- (0,2.3) node[left] {$Y$};
            \draw[thick, color=darkgreen, domain=0.3:2.2, samples=30] plot (\x, {2.0 - 1.8*(\x-1.25)*(\x-1.25)});
            \fill (0.4,0.7) circle (1.5pt); \fill (0.8,1.6) circle (1.5pt);
            \fill (1.25,2.0) circle (1.5pt); \fill (1.7,1.6) circle (1.5pt);
            \fill (2.1,0.7) circle (1.5pt);
            \node[below, font=\scriptsize] at (1.25,-0.2) {$r = 0.0$};
            \node[below, font=\tiny, align=center] at (1.25,-0.6) {Deterministic\\Quadratic Curve};
        \end{scope}
    \end{tikzpicture}
    \caption{Archetypal scatter patterns and their corresponding Pearson correlation values.}
    \label{fig:dm-pearson-scenarios}
\end{figure}

\begin{figure}[htbp]
    \centering
    \includegraphics[width=0.88\textwidth]{iris_correlation_heatmaps.pdf}
    \caption{Correlation heatmaps for the overall Iris dataset and stratified by species.}
    \label{fig:dm-iris-heatmaps}
\end{figure}

\subsection{Causal Fallacy and Confounding Variables}

A cornerstone of data science states that \textbf{correlation does not imply causation}. A high observed $|r_{xy}|$ may arise from three distinct phenomena:
\begin{enumerate}
    \item $X$ is a direct causal driver of $Y$.
    \item $Y$ is a direct causal driver of $X$.
    \item $X$ and $Y$ are simultaneously driven by an unobserved or latent third factor, termed a \textbf{confounding variable}.
\end{enumerate}

\begin{noteblock}{Classic Example: Ice Cream Sales and Drowning Incidents}
During summer months, a high positive correlation ($r > 0.85$) is documented between ice cream sales ($X$) and drowning casualties ($Y$). Inferring that ice cream consumption increases drowning risk represents a classic causal fallacy: both variables are driven by an underlying exogenous confounder—high atmospheric temperature (\textit{hot weather})—which spurs both outdoor swimming and cold dessert purchases.
\end{noteblock}

\subsection{Spearman Rank Correlation Coefficient}

When variables violate bivariate normality assumptions or exhibit non-linear yet strictly monotonic relationships, the \textbf{Spearman Rank Correlation Coefficient ($\rho$)} is preferred.

The method replaces continuous values $x_i$ and $y_i$ with their corresponding ordinal ranks $r(x_i), r(y_i) \in \{1, \dots, n\}$. In the absence of ties, the explicit formulation simplifies to:
\begin{equation}
    \rho = 1 - \frac{6 \sum_{i=1}^n \big(r(x_i) - r(y_i)\big)^2}{n(n^2 - 1)}
\end{equation}

\begin{itemize}[noitemsep]
    \item $\rho = +1$: indicates a perfect monotonic increasing relationship (as $X$ increases, $Y$ strictly increases, even along exponential or logarithmic curves).
    \item $\rho = -1$: indicates a perfect monotonic decreasing relationship.
    \item \textbf{Robustness}: Relying on ranks rather than raw metric differences renders Spearman significantly more resistant to extreme outliers in distribution tails.
\end{itemize}

% ====================================================================
% SECTION 3.3: OUTLIER DETECTION TAXONOMY
% ====================================================================
\section{Outlier Taxonomy and Detection Methodologies}
\label{sec:dm-outlier-detection-taxonomy}

An \textbf{outlier} is defined as an observation that deviates so substantially from the remaining data as to arouse suspicion that it was generated by a different mechanism or contaminated by measurement error.

\subsection{The Dual Nature of Outliers: Noise vs. Primary Target}
\begin{itemize}
    \item \textbf{Outliers as Noise}: The anomaly stems from data entry typos, temporary sensor glitches, or spurious noise. In this role, outliers distort parameter estimation and should be identified, corrected, or removed during data preprocessing.
    \item \textbf{Outliers as the Primary Target (Analysis Goal)}: The explicit mission of the analysis is to uncover rare, high-impact events carrying crucial domain insights:
    \begin{itemize}[noitemsep]
        \item Intrusion detection in computer network security.
        \item Financial fraud detection and anti-money laundering monitoring.
        \item Predictive maintenance through anomalous vibration patterns in industrial turbines.
    \end{itemize}
\end{itemize}

\subsection{Outlier Identification Methodologies}

Outlier identification methodologies can be systematically categorized into univariate and multidimensional approaches:

\begin{figure}[htbp]
    \centering
    \begin{tikzpicture}[
        outlierCard/.style={draw=navy!80!black, fill=navy!6, rounded corners=6pt, thick, font=\footnotesize, align=left, text width=6.0cm, minimum height=2.8cm, inner sep=6pt},
        arrowStyle/.style={-Stealth, thick, draw=navy!85!black}
    ]
        % Root
        \node[rectangle, rounded corners=6pt, draw=purple!85!black, fill=purple!12, very thick, font=\small\bfseries, align=center, text width=10.2cm, inner sep=6pt] (root) at (0, 2.7) {
            OUTLIER DETECTION TAXONOMY\\[2pt]
            \footnotesize \normalfont Criteria for single features and multidimensional spaces
        };

        % Left branch: Univariate
        \node[outlierCard, anchor=north] (uniCol) at (-3.6, 1.1) {
            \textbf{Univariate Methods}\\[4pt]
            $\bullet$ \textbf{Tukey IQR}: Boxplot fences $[L, U]$\\[3pt]
            $\bullet$ \textbf{Z-Score}: Extreme tails $|z| > 3.29$\\[3pt]
            $\bullet$ \textbf{Frequency}: Rare values ($< 0.1\%$)
        };

        % Right branch: Multidimensional
        \node[outlierCard, anchor=north] (multiCol) at (3.6, 1.1) {
            \textbf{Multidimensional Methods}\\[4pt]
            $\bullet$ \textbf{Scatter / SPLOM}: Pairwise outliers\\[3pt]
            $\bullet$ \textbf{PCA}: Residual variance anomalies\\[3pt]
            $\bullet$ \textbf{Clustering / LOF}: Low density
        };

        \coordinate (split) at (0, 1.65);
        \draw[thick, draw=navy!85!black] (root.south) -- (split);
        \draw[arrowStyle] (split) -| (uniCol.north);
        \draw[arrowStyle] (split) -| (multiCol.north);
    \end{tikzpicture}
    \caption{Taxonomy of outlier detection methodologies.}
    \label{fig:dm-outlier-taxonomy}
\end{figure}

\begin{itemize}
    \item \textbf{Categorical Attributes}: An anomalous observation belongs to a rare category occurring with near-zero empirical frequency ($< 0.1\%$).
    \item \textbf{Univariate Numeric Attributes}:
    \begin{itemize}[noitemsep]
        \item \textit{Tukey IQR Rule}: Values outside the interval $[Q_1 - 1.5\,\text{IQR},\, Q_3 + 1.5\,\text{IQR}]$.
        \item \textit{Z-Score Rule}: Standardized values with $|z| > 3.29$ (representing the outermost $0.1\%$ tails of a Gaussian distribution).
    \end{itemize}
    \item \textbf{Multidimensional Spaces}:
    \begin{itemize}[noitemsep]
        \item \textit{Bivariate Scatter Plots}: Visual identification of points detached from dense clusters.
        \item \textit{PCA Projections}: Observations projecting as outliers along secondary principal components.
        \item \textit{Clustering and Local Outlier Factor (LOF)}: Isolating samples residing in low-density neighborhoods relative to their $k$-nearest neighbors.
    \end{itemize}
\end{itemize}

% ====================================================================
% SECTION 3.4: MISSING VALUES
% ====================================================================
\section{Introduction to Missing Values}
\label{sec:dm-missing-values-intro}

A data instance contains \textbf{Missing Values} when no observation is recorded for one or more features (explicitly encoded via symbols like \texttt{NaN}, \texttt{NULL}, \texttt{?} or implicitly hidden under sentinel conventions such as \texttt{999} or \texttt{0}).

\subsection{Root Causes of Missingness}
\begin{enumerate}
    \item \textbf{Voluntary Non-Response}: Respondents decline to disclose sensitive personal data (e.g., annual income, political affiliations, medical history).
    \item \textbf{Conceptual Inapplicability}: The attribute is logically inapplicable for certain subpopulations (e.g., pregnancy counts for male subjects, or employment tenure for school-age minors).
    \item \textbf{Hardware Failures or Dropped Connections}: Telemetry sensors freezing or disconnecting, producing null readings or reporting continuous dummy constants (e.g., a frozen wind sensor persistently reporting $0\,\text{m/s}$).
\end{enumerate}

\subsection{Operational Data Understanding Checklist}

Before proceeding to data cleaning and preparation, practitioners must systematically verify the following diagnostic checklist:

\begin{table}[htbp]
    \centering
    \small
    \begin{tabularx}{\textwidth}{lX}
        \toprule
        \textbf{Diagnostic Check} & \textbf{Operational Verification Procedure} \\
        \midrule
        \textbf{Syntactic Integrity} & Validate data types, string/number formats, and domain constraints. \\
        \textbf{Outlier Inspection} & Inspect extreme values using univariate Boxplots and bivariate Scatter Plots. \\
        \textbf{Missing Value Audit} & Quantify null percentages and detect implicit sentinel values (\texttt{0}, \texttt{-1}, \texttt{999}). \\
        \textbf{Association \& Redundancy} & Compute Pearson and Spearman matrices to identify collinear or redundant features. \\
        \textbf{Distributional Assumptions} & Check symmetry, normality, and variance stability prior to feature transformations. \\
        \bottomrule
    \end{tabularx}
    \caption{Operational diagnostic checklist for completing the Data Understanding phase.}
    \label{tab:dm-du-checklist}
\end{table}
